% Lista03.tex
% !TeX root = L03_RelFinal.tex
\section*{Lista de Exercícios 3}
\addcontentsline{toc}{chapter}{Lista de Exercícios 3}

\noindent \textbf{Data:} 27/09/2026

\section*{Resolução}
\subsection*{Questão 1}
\subsubsection*{a) Modele o problema e encontre as equações discretizadas}

\subsubsection*{Modelagem: Hip\'oteses}

\begin{enumerate}
  \item Condu\c{c}\~ao unidimensional na dire\c{c}\~ao $x$, com temperatura uniforme na se\c{c}\~ao transversal.
  \item Propriedades constantes ($k$, $\rho$, $c$) e se\c{c}\~ao uniforme.
  \item Coeficiente convectivo $h$ e temperatura ambiente $T_\infty$ uniformes
        e constantes.
  \item Topo ($x=L$) isolado; base ($x=0$) mantida em $T_b$.
\end{enumerate}

\subsubsection*{Balan\c{c}o de energia}

Para uma fatia de espessura $dx$ e \'area de se\c{c}\~ao $A$:
\[
  \underbrace{\rho\, c\,A\,dx\,\frac{\partial T}{\partial t}}_{\text{ac\'umulo}} =\underbrace{\left(q_x-q_{x+dx}\right)A}_{\text{condu\c{c}\~ao l\'{\i}quida}} -\underbrace{h\,(\pi D\,dx)\,(T-T_\infty)}_{\text{perda lateral}},
\]
onde $\pi\,D\,dx$ \'e a \'area lateral da fatia. Substituindo $q=-k\,\partial T/\partial x$ (Lei de Fourier) e dividindo toda a equação por por $A\,dx$:
\begin{equation}
  \rho\, c\,\frac{\partial T}{\partial t} =\frac{\partial}{\partial x}\!\left(k\,\frac{\partial T}{\partial x}\right) + S
    \label{eq:edp}
\end{equation}

\vspace{0.1cm}

\begin{equation}
    S=-\frac{h\pi\,D}{A}\left(T-T_\infty\right)
\end{equation}

\vspace{0.1cm}

Sendo o termo $S$ (W/m$^3$) um \emph{sumidouro distribu\'{\i}do}: como o modelo 1D n\~ao possui a superf\'{\i}cie lateral como fronteira, a perda convectiva que nela ocorre \'e diluída no volume de cada fatia.

\subsubsection*{Condi\c{c}\~oes inicial e de contorno}
\[
  T(x,0)=T_\infty ,\qquad
  T(0,t)=T_b ,\qquad
  \left.\frac{\partial T}{\partial x}\right|_{x=L}=0 .
\]

\subsubsection*{Par\^ametros}

\[
\begin{array}{lll}
  T_b=373\ \mathrm{K} & T_\infty=293\ \mathrm{K} & k=10\ \mathrm{W/mK}\\[2pt]
  D=0{,}01\ \mathrm{m} & L=0{,}05\ \mathrm{m} & h=5{,}0\ \mathrm{W/m^2K}\\[2pt]
  \alpha=10^{-6}\ \mathrm{m^2/s} & & 
\end{array}
\]
\subsubsection*{Grandezas derivadas}
\[
  A=\frac{\pi D^{2}}{4}=7{,}854\times10^{-5}\ \mathrm{m^{2}},\
\]
\[
  \rho\, c=\frac{k}{\alpha}=1{,}0\times10^{7}\ \mathrm{J/m^{3}K},
  \qquad
  \frac{h\,\pi D}{A}=2000\ \mathrm{W/m^{3}K}.
\]

\subsubsection*{Discretiza\c{c}\~ao por volumes finitos}

\subsubsection*{Malha}

O dom\'{\i}nio $0\le x\le L$ \'e dividido em $N$ volumes de largura
$\Delta x_P=L/N$, com o n\'o no centro de cada volume
(\emph{cell-centered}). Adota-se a nota\c{c}\~ao: $W$, $P$ e $E$
designam os n\'os oeste, central e leste; $w$ e $e$ as faces do volume;
$\Delta x_W$ e $\Delta x_E$ as dist\^ancias entre n\'os vizinhos. No interior
$\Delta x_W=\Delta x_E=\Delta x_P$; nos volumes de fronteira a dist\^ancia at\'e
o contorno \'e $\Delta x/2$.

\subsubsection*{Integra\c{c}\~ao}

Integra-se a Eq.~\eqref{eq:edp} sobre o volume de controle
e sobre o intervalo de tempo:
\[
  \int_{t}^{t+\Delta t}\!\!\int_{w}^{e}
  \rho c\,\frac{\partial T}{\partial t}\,dx\,dt
  =
  \int_{t}^{t+\Delta t}\!\!\int_{w}^{e}
  \left[\frac{\partial}{\partial x}\!\left(k\frac{\partial T}{\partial x}\right)
  +S\right]dx\,dt .
\]

Resultando em:
\[
  \rho\, c\,\,\Delta x_P\,\left(\frac{T_P-T_P^{0}}{\Delta t}\right)
  =\frac{k}{\Delta x_E}\left(T_E-T_P\right)
  -\frac{k}{\Delta x_W}\left(T_P-T_W\right)
  -\frac{h\pi D\,\Delta x_P}{A}\left(T_P-T_\infty\right),
\]
em que o sobrescrito $0$ indica o tempo anterior.

\subsubsection*{Equa\c{c}\~ao discretizada -- volumes internos ($2\le i\le N-1$)}

Agrupando por inc\'ognita, chega-se \`a forma can\^onica
\begin{equation}
  \boxed{\;A_P\,T_P=A_W\,T_W+A_E\,T_E+B\;}
  \label{eq:geral}
\end{equation}
com
\[
\begin{aligned}
  A_W&=\frac{k}{\Delta x_W}=\frac{k}{\Delta x_P}, &
  A_E&=\frac{k}{\Delta x_E}=\frac{k}{\Delta x_P}, \\[4pt]
  A_P^{0}&=\frac{\rho c\,\Delta x_P}{\Delta t}, &
  A_P&=A_W+A_E+A_P^{0}+\dfrac{h\pi D}{A}\,\Delta x_P, \\[4pt]
  B&=A_P^{0}\,T_P^{0}+\dfrac{h\pi D}{A}\,\Delta x_P\,T_\infty . &&
\end{aligned}
\]

Escrevendo o termo fonte linearizado como $S=S_C+S_P\,T_P$, com
\[
  S_C=\frac{hP}{A}\,T_\infty ,\qquad S_P=-\frac{hP}{A}\le 0 ,
\]
a parcela constante integra $B$ e a parcela proporcional a $T_P$ soma-se a
$A_P$. A condi\c{c}\~ao $S_P\le 0$ garante
$A_P\ge A_W+A_E$, isto \'e, domin\^ancia diagonal da matriz.

\subsubsection*{Volume de fronteira 1 -- base com temperatura prescrita}

A face oeste do primeiro volume coincide com a base, onde $T=T_b$ \'e
conhecido; logo $\Delta x_W=\Delta x_P/2$ e o produto $A_W T_b$ passa ao termo independente:
\begin{equation}
  A_P\,T_1=A_E\,T_2+B_1 ,
\end{equation}
\[
\begin{aligned}
  A_E&=\frac{k}{\Delta x_P}, &
  A_W&=\frac{2k}{\Delta x_P}, \\[4pt]
  A_P&=A_E+A_W+A_P^{0}+\dfrac{h\pi D\,\Delta x_P}{A}, &
  B_1&=A_W\,T_b+A_P^{0}\,T_1^{0}+\dfrac{\pi D\,h\,\Delta x_P}{A}\,T_\infty .
\end{aligned}
\]

\subsubsection*{Volume de fronteira $N$ -- topo isolado}

Na face leste o fluxo \'e nulo, de modo que o termo correspondente n\~ao
comparece no balan\c{c}o:
\begin{equation}
  A_P\,T_N=A_W\,T_{N-1}+B_N ,
\end{equation}
\[
\begin{aligned}
  A_W&=\frac{k}{\Delta x_P}, &
  A_E&=0, \\[4pt]
  A_P&=A_W+A_P^{0}+\dfrac{\pi D\,h\,\Delta x_P}{A}, &
  B_N&=A_P^{0}\,T_N^{0}+\dfrac{\pi D\,h\,\Delta x_P}{A}\,T_\infty .
\end{aligned}
\]

\subsubsection*{Sistema resultante}

As $N$ equa\c{c}\~oes formam, a cada passo de tempo, um sistema linear
tridiagonal
\[
  -A_W\,T_{m-1}+A_P\,T_m-A_E\,T_{m+1}=B_m ,\qquad m=1,\dots,N,
\]
com $A_W=0$ em $m=1$ e $A_E=0$ em $m=N$, resolvido diretamente pelo algoritmo
de Thomas (TDMA). Como $k$, $\Delta x_P$ e $\Delta t$ s\~ao constantes, os
coeficientes $A_W$, $A_E$ e $A_P$ s\~ao calculados uma \'unica vez; apenas o
termo fonte $B_m$ \' e atualizado a cada avan\c{c}o no tempo, por conter
$T_i^{0}$.

\vspace{0.5cm}

\subsubsection*{b) Implemente a solução utilizando o TDMA}

TDMA:

O método admite que cada incógnita se escreva em função apenas da seguinte,
\begin{equation}
  T_m = P_m\,T_{m+1} + Q_m ,
\end{equation}
e determina os coeficientes de recorrência por substituição na equação do
volume $m$:
\begin{equation}
  P_m=\frac{A_E}{A_P-A_W P_{m-1}},\qquad
  Q_m=\frac{B_m+A_W Q_{m-1}}{A_P-A_W P_{m-1}},\qquad P_0=Q_0=0 .
\end{equation}

A varredura para frente ($m=1,\dots,N$) monta as $N$ relações sem resolver
nenhuma temperatura. A varredura para trás inicia no último volume, onde
$A_E=0$ implica $P_N=0$ e portanto $T_N=Q_N$; os demais valores seguem por
substituição regressiva, $T_m=P_mT_{m+1}+Q_m$, com $m=N-1,\dots,1$. Os zeros
$P_0=0$ e $P_N=0$ são as duas fronteiras do domínio fechando o algoritmo.

Como o problema é linear ($k$ constante e termo fonte linear em $T$), o TDMA é
um método \emph{direto}: cada passo de tempo é resolvido exatamente, sem
iteração interna.

\subsubsection*{Marcha no tempo}

Os coeficientes $A_W$, $A_E$ e $A_P$ são constantes (malha fixa, $\Delta t$
fixo, propriedades constantes) e por isso montados uma única vez, fora do
laço. Apenas $B_m$ é atualizado a cada avanço, por conter $T_m^{0}$.

O procedimento é:

\begin{enumerate}
  \item Inicializar $T_m = T_\infty$ em todos os volumes;
  \item Montar $B_m$ com o campo do tempo anterior, acrescentando
        $A_W T_b$ em $B_1$;
  \item Resolver o sistema pelo TDMA, obtendo o campo no tempo novo;
  \item Avaliar $\max_i|T_m^{i+1}-T_m^{i}|$; se inferior a $\varepsilon$,
        encerrar; caso contrário, fazer $T^{0}\leftarrow T$ e retornar ao
        passo~2.
\end{enumerate}

Ressalte-se que o critério $\varepsilon$ aplica-se à \emph{marcha temporal} --
isto é, à aproximação ao regime permanente -- e não à solução do sistema
linear, que é exata a cada passo.

\subsubsection*{Verificação}

Para $N=10$, $\Delta t=1$~s e $\varepsilon=10^{-6}$~K, a convergência ocorreu
em 9723 passos.

\vspace{0.5cm}

\subsubsection*{c) Resolva o problema variando o número de volumes elementares}

\begin{table}[H]
\centering
\caption{Influência do refino da malha ($\Delta t=1$~s,
$\varepsilon=10^{-6}$~K). O erro é o desvio máximo nodal em relação à solução
analítica do item~(e).}
\begin{tabular}{crrrr}
\hline
$N$ & $\Delta x_P$ (m) & $|e|_{\max}$ (K)\\
\hline
3  & 0,01667 & 0,49408\\
5  & 0,01000 & 0,18668\\
10 & 0,00500 & 0,04838\\
20 & 0,00250 & 0,01232\\
40 & 0,00125 & 0,00312\\
80 & 0,00063 & 0,00137\\
\hline
\end{tabular}
\end{table}

Pode-se perceber que $|e|_{max}$ diminui a medida que o número de elementos de malha aumenta.

\vspace{0.5cm}

\subsubsection*{d) Resolva o problema variando o critério de convergência}

\begin{table}[H]
\centering
\caption{Influência do critério de parada ($N=20$, $\Delta t=1$~s).}
\begin{tabular}{crrrrr}
\hline
$\varepsilon$ (K) & passos & $t$ (s) & $|e|_{\max}$ (K)\\
\hline
$10^{-1}$ &   188 &   188 & 61,861\\
$10^{-2}$ &  1946 &  1946 & 8,432\\
$10^{-4}$ &  5830 &  5830 & 0,0927\\
$10^{-6}$ &  9714 &  9714 & 0,01232\\
$10^{-8}$ & 13597 & 13597 & 0,01229\\
\hline
\end{tabular}
\end{table}

Critérios frouxos ($\varepsilon\ge10^{-2}$~K) interrompem a marcha antes do
regime permanente; para $\varepsilon=10^{-1}$~K a ponta da aleta permanece
praticamente na temperatura inicial. No extremo oposto observa-se saturação:
de $10^{-6}$ para $10^{-8}$~K o custo cresce $40\%$ e o erro passa de
$0{,}01232$ para $0{,}01229$~K. 
\vspace{0.5cm}

\subsubsection*{e) Compare a solução de regime permanente com a solução analítica}

Em regime permanente a Eq.~\eqref{eq:edp} reduz-se a $\theta''-m^2\theta=0$,
com $\theta=T-T_\infty$ e $m=\sqrt{h\pi D/kA}=\sqrt{4h/kD}$. Com
$\theta(0)=\theta_b$ e $\theta'(L)=0$, a solução clássica para aleta de ponta
adiabática é
\begin{equation}
  \frac{T(x)-T_\infty}{T_b-T_\infty}
  =\frac{\cosh\!\left[m(L-x)\right]}{\cosh(mL)},
  \qquad
  \dot Q_{\text{base}}=\sqrt{h\pi D\,kA}\,\left(T_b-T_\infty\right)\tanh(mL).
\end{equation}

Com os dados do problema:
\[
  m=\sqrt{\frac{4\times5{,}0}{10\times0{,}01}}=14{,}142\ \mathrm{m^{-1}},
  \qquad mL=0{,}7071,\qquad \cosh(mL)=1{,}2606 .
\]

\begin{table}[H]
\centering
\caption{Campo de temperatura em regime permanente ($N=10$,
$\Delta t=1$~s, $\varepsilon=10^{-6}$~K). Nós em
$x_m=(m-\tfrac{1}{2})\Delta x_P$.}
\begin{tabular}{crrrr}
\hline
$m$ & $x$ (m) & $T_{\text{num}}$ (K) & $T_{\text{teórico}}$ (K) & erro (K)\\
\hline
1  & 0,0025 & 371,279 & 371,328 & $-0{,}0484$\\
2  & 0,0075 & 368,229 & 368,274 & $-0{,}0455$\\
3  & 0,0125 & 365,555 & 365,598 & $-0{,}0430$\\
4  & 0,0175 & 363,243 & 363,284 & $-0{,}0409$\\
5  & 0,0225 & 361,283 & 361,322 & $-0{,}0391$\\
6  & 0,0275 & 359,665 & 359,702 & $-0{,}0377$\\
7  & 0,0325 & 358,379 & 358,416 & $-0{,}0365$\\
8  & 0,0375 & 357,421 & 357,456 & $-0{,}0357$\\
9  & 0,0425 & 356,784 & 356,820 & $-0{,}0351$\\
10 & 0,0475 & 356,467 & 356,502 & $-0{,}0348$\\
\hline
\end{tabular}
\end{table}

\subsection*{Questão 2}
\subsubsection*{a) Por que o problema proposto não tem solução? Explique a razão.}
 
O problema não tem solução, porque, simplesmente, o balanço de energia na
placa plana não fecha.
 
Resolvendo a EDO:
\[
  \frac{d^{2}T}{dx^{2}}+\frac{q'''}{k}=0
  \quad\therefore\quad
  \int \frac{d^{2}T}{dx^{2}}\,dx = -\int \frac{q'''}{k}\,dx
\]
\[
  \frac{dT}{dx}=-\frac{q'''}{k}\,x+C
  \quad\therefore\quad
  \int \frac{dT}{dx}\,dx = \int\left(-\frac{q'''}{k}\,x+C\right)dx
\]

\[T(x)=-\frac{q'''}{2k}\,x^{2}+C_1\,x+C_2\]
 
Aplicando as condições de contorno:
\[
  -k\left.\frac{dT}{dx}\right|_{x=0}=10\ \mathrm{W/m^{2}}
\]
\[
  -k\left.\frac{dT}{dx}\right|_{x=L}=21\ \mathrm{W/m^{2}}
\]
 
Sejam:
\[
  q'''=7\ \mathrm{W/m^{3}}\ ;\quad k=1\ \mathrm{W/(m\,K)}\ ;\quad L=3\ \mathrm{m}
\]
 
Primeira condição de contorno:
\[
  -\frac{10}{1}=-\frac{7}{1}\times 0+C
  \quad\therefore\quad
  C=-10\ \mathrm{K/m}
\]
 
Segunda condição de contorno:
\[
  -\frac{21}{1}=-\frac{7}{1}\times 3+C
  \quad\therefore\quad
  C=0\ \mathrm{K/m}
\]
 
Como, para atender as condições de contorno impostas à EDO, devemos ter dois
valores diferentes para a constante $C$, o problema, da forma com que foi
proposto, não possui solução.

\subsubsection*{b) Altere a condição de contorno que deixa a placa, e considere
que o fluxo de calor que sai é igual a 31 W/m$^2$. Por que não é possível
encontrar a solução usando TDMA?}

Com $q(L)=31\ \mathrm{W/m^2}$ o balanço global passa a fechar, pois, as condições passam a ser:

\[
  -k\left.\frac{dT}{dx}\right|_{x=0}=10\ \mathrm{W/m^{2}}
\]
\[
  -k\left.\frac{dT}{dx}\right|_{x=L}=31\ \mathrm{W/m^{2}}
\]

E os cálculos de $C$ são: 

Primeira condição de contorno:
\[
  -\frac{10}{1}=-\frac{7}{1}\times 0+C
  \quad\therefore\quad
  C=-10\ \mathrm{K/m}
\]
 
Segunda condição de contorno:
\[
  -\frac{31}{1}=-\frac{7}{1}\times 3+C
  \quad\therefore\quad
  C=-10\ \mathrm{K/m}
\]

Como $C$ é igual para ambas as condições de contorno, podemos dizer que o balanço energético do problema atende à Primeira Lei da Termodinâmica.

Ainda assim o TDMA não resolve o problema. Para evidenciar a razão, aplica-se
o algoritmo a uma malha de $N=3$ volumes, para a qual
$\Delta x=1\ \mathrm{m}$ e $k/\Delta x=1\ \mathrm{W/m^2K}$.

\subsubsection*{Coeficientes}

Os dois contornos prescrevem fluxo, de modo que nenhum deles contribui para
$A_P$: um fluxo imposto não depende da temperatura do nó e vai integralmente
para o termo fonte.

\begin{table}[H]
\centering
\caption{Coeficientes para $N=3$, com fluxo prescrito nas duas faces.}
\begin{tabular}{ccccc}
\hline
$m$ & $A_W$ & $A_E$ & $A_P$ & $B$\\
\hline
1 & 0 & 1 & 1 & $q'''\Delta x + q(0) = 7+10 = 17$\\
2 & 1 & 1 & 2 & $q'''\Delta x = 7$\\
3 & 1 & 0 & 1 & $q'''\Delta x - q(L) = 7-31 = -24$\\
\hline
\end{tabular}
\end{table}

\subsubsection*{Varredura para frente}

\[
\begin{aligned}
  m=1:&\quad P_1=\frac{A_E}{A_P}=\frac{1}{1}=1, &
              &Q_1=\frac{B_1}{A_P}=\frac{17}{1}=17\\[4pt]
  m=2:&\quad A_P-A_WP_1=2-1(1)=1, &
              &P_2=\frac{1}{1}=1,\qquad
               Q_2=\frac{B_2+A_WQ_1}{1}=\frac{7+17}{1}=24\\[4pt]
  m=3:&\quad A_P-A_WP_2=1-1(1)=0 &&
\end{aligned}
\]

O denominador se anula no último volume e a recorrência é interrompida: não é
possível calcular $P_3$ nem $Q_3$. O numerador correspondente também é nulo,
\[
  B_3+A_WQ_2=-24+24=0 ,
\]
de modo que a divisão resulta na indeterminação $0/0$. Dessa forma, não é possível seguir com o método para solução das equações.

Isso ocorrre, pois a diagonal principal apenas iguala a soma dos demais termos da matriz, sem superá-la,
o que leva a $P_m=1$ e ao denominador nulo em $m=N$.

\subsubsection*{c) Fixe o valor da temperatura da fronteira à direita em
10\textdegree C. Por que agora foi possível achar a solução usando TDMA?}

Com temperatura prescrita à direita, o contorno passa a estar a meia largura
de volume do nó~3, o que introduz o  termo $2k/\Delta x$, ligando esse
nó a um valor conhecido de temperatura (o valor prescrito). Essa parcela soma-se a $A_P$, e o produto entre esse termo e o novo valor prescrito de temperatura vai para $B$:

\begin{table}[H]
\centering
\caption{Coeficientes para $N=3$, com $T=10$\textdegree C na face direita.}
\begin{tabular}{ccccc}
\hline
$m$ & $A_W$ & $A_E$ & $A_P$ & $B$\\
\hline
1 & 0 & 1 & 1 & $7+10 = 17$\\
2 & 1 & 1 & 2 & $7$\\
3 & 1 & 0 & $1+2 = 3$ & $7 + 2(10) = 27$\\
\hline
\end{tabular}
\end{table}

No volume~3 a diagonal passa a superar a soma dos demais termos
($A_P=3 > A_W+A_E=1$), o que restabelece a dominância diagonal estrita.
Refazendo a varredura para frente:
\[
\begin{aligned}
  m=1:&\quad P_1=\frac{1}{1}=1, &&Q_1=\frac{17}{1}=17\\[4pt]
  m=2:&\quad A_P-A_WP_1=2-1(1)=1, &&P_2=1,\qquad Q_2=\frac{7+17}{1}=24\\[4pt]
  m=3:&\quad A_P-A_WP_2=3-1(1)=2, &&P_3=0,\qquad Q_3=\frac{27+24}{2}=25{,}5
\end{aligned}
\]

O denominador não se anula mais, e a varredura para trás fornece a solução:
\[
  T_3=Q_3=25{,}5\ \text{\textdegree C},\qquad
  T_2=P_2T_3+Q_2=49{,}5\ \text{\textdegree C},\qquad
  T_1=P_1T_2+Q_1=66{,}5\ \text{\textdegree C}.
\]

\subsubsection*{d) Compare as soluções do item c) para os diferentes métodos
empregados}

Os quatro métodos foram aplicados à mesma malha de $N=5$ volumes, com
tolerância $\varepsilon=10^{-8}$~K para os métodos iterativos.

\begin{table}[H]
\centering
\caption{Temperaturas nodais obtidas pelos quatro métodos
($N=5$, $\varepsilon=10^{-8}$).}
\begin{tabular}{ccccccc}
\hline
$m$ & $x$ (m) & Jacobi & Gauss-Seidel & SOR & TDMA & Analítica\\
\hline
1 & 0,300 & 68,500 & 68,500 & 68,500 & 68,500 & 68,185\\
2 & 0,900 & 59,980 & 59,980 & 59,980 & 59,980 & 59,665\\
3 & 1,500 & 48,940 & 48,940 & 48,940 & 48,940 & 48,625\\
4 & 2,100 & 35,380 & 35,380 & 35,380 & 35,380 & 35,065\\
5 & 2,700 & 19,300 & 19,300 & 19,300 & 19,300 & 18,985\\
\hline
\end{tabular}
\end{table}

\begin{table}[H]
\centering
\caption{Custo de cada método ($N=5$, $\varepsilon=10^{-8}$).}
\begin{tabular}{lc}
\hline
Método & iterações\\
\hline
Jacobi               & 330\\
Gauss-Seidel         & 168\\
SOR ($\omega=1{,}2$) & 111\\
SOR ($\omega=1{,}5$) &  36\\
SOR ($\omega=1{,}8$) & 100\\
TDMA                 & ---\\
\hline
\end{tabular}
\end{table}

Todos os métodos convergem para a mesma solução, pois resolvem o mesmo sistema
linear; a diferença está unicamente no custo computacional. O Gauss-Seidel
exige cerca de metade das iterações do Jacobi, por utilizar os valores já
atualizados dentro da própria varredura. O SOR reduz ainda mais esse número,
desde que $\omega$ seja adequado: com $\omega=1{,}5$ foram necessárias apenas
36 iterações, enquanto $\omega=1{,}2$ e $\omega=1{,}8$ elevaram o total para
111 e 100, respectivamente, evidenciando a existência de um valor ótimo.

O TDMA não itera. Por ser um método direto, resolve o sistema tridiagonal em
duas varreduras, fornecendo a solução exata do sistema discreto e dispensando
critério de convergência.

O desvio em relação à solução analítica é constante e igual a
0,315~K em todos os nós, comum aos quatro métodos por se tratar de
erro de discretização, e não de solução do sistema.